% !TITLE 诗经·周南·关雎
% !AUTHOR 无名氏
% !DYNASTY 先秦
% !GENRE 四言诗
% !ORDER 1
% !SOURCE 诗经·周南

\begin{poem}
关关雎鸠\textsuperscript{1}，在河之洲。窈窕淑女，君子好逑\textsuperscript{2}。

参差荇菜\textsuperscript{3}，左右流之。窈窕淑女，寤寐求之。
求之不得，寤寐思服。悠哉悠哉，辗转反侧\textsuperscript{4}。

参差荇菜，左右采之。窈窕淑女，琴瑟友之。
参差荇菜，左右芼之。窈窕淑女，钟鼓乐之。
\end{poem}

\begin{annotations}
\notetext{1}{关关：水鸟和鸣的声音。雎鸠：一种水鸟，雌雄情意专一，古人常用以比喻忠贞的爱情。}
\notetext{2}{好逑：好的配偶。逑，匹配之意。君子好逑指君子美好的配偶。}
\notetext{3}{荇菜：一种水生植物，嫩叶可食。诗中女子采荇菜的动作，暗示了勤劳贤淑的品德。}
\notetext{4}{辗转反侧：翻来覆去睡不着的样子，形容思念之深。}
\end{annotations}

\begin{appreciation}
《诗经》三百零五篇，以《关雎》开卷。这不是偶然的。

孔子删诗，定《诗》三百，将这样一首爱情诗置于卷首，这件事本身就值得玩味。他不是不知道礼教的重量，但他仍然选择了它。因为在孔子看来，男女之情不是礼教的对立面，而恰恰是礼乐文明的起点。夫妇之道，人伦之始——爱一个人，发乎情，止乎礼，这本身就是儒家最看重的修养功夫。

《关雎》的好，在于它写出了爱情里的节制。君子爱慕淑女，从"寤寐求之"到"辗转反侧"，思念到了睡不着觉的程度，但全诗没有一个字越界。他没有用粗暴的方式去表白，更没有任何轻浮的举止，而是想着"琴瑟友之""钟鼓乐之"——用音乐去亲近她，用礼仪去迎娶她。这大概就是孔子说的"哀而不伤"——情感可以很深，但不至于沉溺。也是后世所讲的"温柔敦厚"——心里翻江倒海，面上依旧温润如玉。它有爱，但不滥，不淫。

这种克制的美，到了明代汤显祖那里，成了一声惊雷。《牡丹亭》里的杜丽娘读到"窈窕淑女，君子好逑"，一夜之间从循规蹈矩的闺秀变成了为情生死的痴人。她不是被这首诗教坏的——恰恰相反，是这首诗让她第一次意识到，男女之情可以这样自然、这样正当、这样被圣贤郑重地对待。一首三千年前的诗，隔了重重帷幕，仍然能在一个少女心里掀起这样滔天的波澜，这便是诗歌真正的力量。它不靠说教，只是让你看见：爱，原来可以是这样子的。
\end{appreciation}
